\chapter{Auswertung}
\label{ch:Auswertung}
% Liste der genutzer Formeln für die Fehlerrechnung
\section*{Fehlerrechnung}
Für die statistische Auswertung von $n$ Messwerten $x_i$ werden folgende Größen definiert \cite{errorSkript25}:
\begin{align}
    \bar{x} &= \frac{1}{n} \sum_{i=1}^{n} x_i \vphantom{\sqrt{\sum_i^n}^2} && \text{\textcolor{gray}{Arithmetisches Mittel}} \label{eq:arithmetisches_mittel} \\
    \sigma^2 &= \frac{1}{n-1} \sum_{i=1}^{n} (x_i - \bar{x})^2 \vphantom{\sqrt{\sum_i^n}^2} && \text{\textcolor{gray}{Varianz}} \label{eq:Varianz} \\
    \sigma &= \sqrt{\frac{1}{n-1} \sum_{i=1}^{n} (x_i - \bar{x})^2} \vphantom{\sqrt{\sum_i^n}^2} && \text{\textcolor{gray}{Standardabweichung}} \label{eq:standardabweichung} \\
    \Delta \bar{x} &= \frac{\sigma}{\sqrt{n}} = \sqrt{\frac{1}{n(n-1)} \sum_{i=1}^n(\bar x - x_i)^2} \vphantom{\sqrt{\sum_i^n}^2} && \text{\textcolor{gray}{Fehler des Mittelwerts}} \label{eq:fehler_mittelwert} \\
    \Delta f &= \sqrt{\left(\frac{\partial f}{\partial x} \Delta x\right)^2 + \left(\frac{\partial f}{\partial y} \Delta y\right)^2} \vphantom{\sqrt{\sum_i^n}^2} && \text{\textcolor{gray}{Gauß’sches Fehlerfortpflanzungsgesetz für $f(x,y)$}} \label{eq:gauss_fehlfortpflanzung} \\
    \Delta f &= \sqrt{(\Delta x)^2 + (\Delta y)^2} \vphantom{\sqrt{\sum_i^n}^2} && \text{\textcolor{gray}{Fehler für $f = x + y$}} \label{eq:fehler_summe} \\
    \Delta f &= |a| \Delta x \vphantom{\sqrt{\sum_i^n}^2} && \text{\textcolor{gray}{Fehler für $f = ax$}} \label{eq:fehler_proportional} \\
    \frac{\Delta f}{|f|} &= \sqrt{\left(\frac{\Delta x}{x}\right)^2 + \left(\frac{\Delta y}{y}\right)^2} \vphantom{\sqrt{\sum_i^n}^2} && \text{\textcolor{gray}{relativer Fehler für $f = xy$ oder $f = x/y$}} \label{eq:relativer_fehler} \\
    \sigma &= \frac{|a_{lit} - a_{gem}|}{\sqrt{\Delta a_{lit}^2 + \Delta a_{gem}^2}} \vphantom{\sqrt{\sum_i^n}^2} && \text{\textcolor{gray}{Berechnung der signifikanten Abweichung}} \label{eq:signifikante_abweichung}
\end{align}

\newpage

\section{Bestimmung der Halbwertszeit von Silber-Isotropen}
\label{sec:A1}
\subsection{Untergrundsstrahlung bei der Silbermessung}
In dieser Sektion soll der Zerfall des Silbers betrachtet werden. Dafür wurde zunächst ohne Probe die Untergrundsstrahlung bestimmt. Dazu wurde für die Torzeit 10s und für die Messzeit 500s engestellt. Es wird das \hyperref[eq:arithmetisches_mittel]{arithmetische Mittel (\ref*{eq:arithmetisches_mittel})} und seine Ungenauigkeit über den \hyperref[eq:fehler_mittelwert]{Fehler des Mittelwertes (\ref*{eq:fehler_mittelwert})} bestimmt. 
Zusätzlich wird die Messung mit einem Faktor $n$ multiplizietr, wobei $n$ die Anzahl der Messreihen ist. Für diese Silbermessung gilt $n=4$ 

\eq{avg_unter_ag}{
    \overline{N_\text{UG,Ag}} = () s^{-1}
}

\subsection{Bestimmung der Halbwertszeit}
Die vier Messreihen werden aufsummiert umd eine Gesamtmessreihe zu erhalten. Diese wird gegend die Zeit geplottet. Durch die Messdaten wird eine Regression der Form
\eq{reg_ag}{
   f(x) = A_1 e^{-\lambda_1 x} + A_2 e^{-\lambda_2 x} + y_0 \, . 
} 
gelegt. Dass es eine Summation aus zwei Exponentialfunktionen entspricht liegt daran, dass zwei Isotrope von Silber natürlich vorkommen. 
Der Index 1 soll das Isotrop $^{110}\mathrm{Ag}$ beschreiben und der Index 2 das Isotrop $^{108}\mathrm{Ag}$. Der Term $y_0$ beschreibt die Untergrundsstrahlung.
Die Messdaten mit der Regression sind der \cabb{plots/Regression_der_Messdaten_1a} zu entnehmen. 
\img{plots/Regression_der_Messdaten_1a}{Bestimmte Regressionen für die Silbermessung.}

Es werden zwei weitere Fits erstellt, welche dazu dienen, eine präzisere Ungenauigkeit der Untergrundsmessung zu bestimmen. Dazu wird der Untergrund über- und unterschätzt, so ergeben sich ein Maximal- und ein Minimalwert. Aus denen wird das \hyperref[eq:arithmetisches_mittel]{arithmetische Mittel (\ref*{eq:arithmetisches_mittel})} gebildet, dies entsprcht der Ungenauigkeit der gemessenen Untergrundsstrahlung.
Es gilt also die Differenzen $|\lambda - \lambda_\text{min}|$ und $|\lambda - \lambda_\text{max}|$ zu bestimmen. Aus diesen wird das \hyperref[eq:arithmetisches_mittel]{arithmetische Mittel (\ref*{eq:arithmetisches_mittel})} gebildet. 
Nach \ceq{fehler_summe} können nun die Ungenauigkeit der Untergrundsmessung und dem \hyperref[eq:fehler_mittelwert]{Fehler des Mittelwertes (\ref*{eq:fehler_mittelwert})} verrechnet werden. Dies ist der Gesamtfehler der Zerfallskonstanten $\Delta \lambda$.

Die Fitparameter beider Regressionen und die Regression selsbt sind der \cabb{plots/Regression_der_Messdaten_1b} zu entnehmen.
% \com{Juan hat hier eine advanst Fehler bestimmung gemacht --> Nachvollziehen und klauen}
\img{plots/Regression_der_Messdaten_1b}{Bestimmte Regressionen mit Unter- und Überschätzung für die Silbermessung.}



Aus der Zefallskosntanten lassen sich die Halbwertszeit $T_{1/2}$ nach \ceq{halbwertszeit} und die Lebensdauer $\tau$ nach \ceq{lebensdauer} bestimmen. Die Ergebnisse sind der \ctab{AT1} zu entnehmen.
\begin{table}[h!]
    \centering
    \caption{Messergebnisse der Silbermessung mit \hyperref[eq:standardabweichung]{Standardabweichung (\ref*{eq:standardabweichung})} zum Literaturwert nach \cabb{Nuklioide}.}
    \label{tab:AT1}
    \begin{tabular}{L | C | C | C | C | C}
    \toprule
    \text{Größe} &
    \text{Gesamt} &
    \overline{N_\text{UG,Ag}} &
    \overline{N_\text{UG,Ag}} + \Delta \overline{N_\text{UG,Ag}} &
    \overline{N_\text{UG,Ag}} - \Delta \overline{N_\text{UG,Ag}} &
    \text{Std.Abw.} \sigma \\ 
    \midrule
    \lambda_1 [s^-1] &  &  &  &  &  \\
    \lambda_2 [s^-1] &  &  &  &  &  \\
    T_{1/2}^1 [s]    &  &  &  &  &  \\
    T_{1/2}^2 [s]    &  &  &  &  &  \\
    \tau_1    [s]    &  &  &  &  &  \\
    \tau_2    [s]    &  &  &  &  &  \\ 
    \bottomrule
    \end{tabular}
\end{table}

\section{Bestimmung der Halbwertszeit von Indium-Isotropen}
\label{sec:A2}
Analog zur Auswerung des Silbers in der \csec{A1} soll hier die Indiumprobe untersucht werden. Hier ist nur ein Isotrop relevant, denn $^{116m}In$ hat eine Halbwertszeit von $T_{1/2} = 14$s, wohingegen $^{116}In$ eine Halbwertszeit von $T_{1/2} = 54$min hat. Da dir Torzeit auf 10s gestellt ist, muss der erste Messwert herausgenommen werden, da deiser den Zerfall von $^{116}In$ im großen Maße mit gemessen hat. 
Die Regression für die Indiummessung ergitb nach \ceq{aktivitaet_decay} sich zu
\eq{reg_in}{
   f(x) = A_0 e^{-\lambda_1 x} + y_0 \, . 
} 
Dabei sind die Variablen wie in \ceq{reg_ag} zu interpretieren. Auch hier muss der Untergrund betrachtet werden. Da dieses Mal die Torzeit auf 120s gestellt wurde, muss das \hyperref[eq:arithmetisches_mittel]{arithmetische Mittel (\ref*{eq:arithmetisches_mittel})} der Untergrundsmessung mit 12 multipilziert werden.
Es ergibt sich:
\eq{avg_unter_in}{
    \overline{N_\text{UG,In}} = () s^{-1}
}
Die Messdaten, sowie die Regression sind der \cabb{plots/Regression_der_Messdaten_2a} zu entnehmen. 
\img{plots/Regression_der_Messdaten_2a}{Bestimmte Regressionen für die Indiummessung.}

Auch hier soll der Untergrund genauer bestimmt werden, es werden wieder eine Maximal- und Minimalabschätzung unternommen. Die Regressionen und deren Fitparameter sind der \cabb{plots/Regression_der_Messdaten_2b} zu entnehmen.


\img{plots/Regression_der_Messdaten_2b}{Bestimmte Regressionen mit Unter- und Überschätzung für die Indiummessung.}

\begin{table}[h!]
    \centering
    \caption{Messergebnisse der Indiummessung mit \hyperref[eq:standardabweichung]{Standardabweichung (\ref*{eq:standardabweichung})} zum Literaturwert nach \cabb{Nuklioide}.}
    \label{tab:AT2}
    \begin{tabular}{L | C | C | C | C | C}
    \toprule
    \text{Größe} &
    \text{Gesamt} &
    \overline{N_\text{UG,In}} &
    \overline{N_\text{UG,In}} + \Delta \overline{N_\text{UG,In}} &
    \overline{N_\text{UG,In}} - \Delta \overline{N_\text{UG,In}} &
    \text{Std.Abw.} \sigma \\ 
    \midrule
    \lambda [s^-1] &  &  &  &  &  \\
    T_{1/2} [s]    &  &  &  &  &  \\
    \tau    [s]    &  &  &  &  &  \\
    \bottomrule
    \end{tabular}
\end{table}

\section{Bestimmung der Güte der Regressionen}
\label{sec:A3}
Zuletzt sollen die bestimmten Regressionen auf ihre Güte geprüft werden. Dazu wird zu nächst das Chi-Quadrat ($\chi^2$) bestimmt. Zudem wird dieses via Division mit den Freiheitsgraden dividert. Ist das $\chi_\text{red}^2 \approx 1$\footnote{Ein Wert von $=1$ wäre ein optimaler Fit.}, ist der Fit als sinnvoll zu erachten.
Die Freiheitsgrade können via
\eq{DoF}{
    \text{DoF\footnote{DoF = Degrees of Freedom}} = n_\text{mess} - n_\text{par}
}
bestimmt werden.
Zuletzt kann die Güte einer Regression noch über die Wahrscheinlichkeit eines Fits bestimmt werden. Diese Methode nimmt die bestimmten Regressionsparameter als wahr an und verteilt die Messwerte normal um den Mittelwert. Die Übereinstimmung der theoretischen und experimentellen Werte wird dann bestimmt.
Die bestimmten Werte sind der \ctab{AT3} zu entnehmen.

\begin{table}[h!]
    \centering
    \caption{Messergebnisse der Indiummessung mit \hyperref[eq:standardabweichung]{Standardabweichung (\ref*{eq:standardabweichung})} zum Literaturwert nach \cabb{Nuklioide}.}
    \label{tab:AT3}
    \begin{tabular}{L | C | C}
    \toprule
    \text{Größe} &
    \text{Silber} &
    \text{Indium} \\
    \midrule
    \chi^2                      &  &  \\
    \chi^2_\text{red}           &  &  \\
    \text{Wahrscheinlichkeit}   &  &  \\
    \bottomrule
    \end{tabular}
\end{table}